% =====================================================================
%  Resilient Collaborative Perception for Multi-Robot Navigation under
%  Sensor Noise and Byzantine False-Obstacle Attacks
%  Target: Robotics and Autonomous Systems (Elsevier) - elsarticle class
%  Build: upload this folder to Overleaf (template: "Elsevier article")
% =====================================================================
\documentclass[preprint,12pt]{elsarticle}
% For final double-column proof view use: \documentclass[final,5p,times,twocolumn]{elsarticle}

\usepackage{amsmath,amssymb}
\usepackage{graphicx}
\usepackage{booktabs}      % clean tables (no vertical rules, per Elsevier)
\usepackage{multirow}
\usepackage[ruled,vlined,linesnumbered]{algorithm2e}
\usepackage{xcolor}
\usepackage[hidelinks]{hyperref}
\newcommand{\todo}[1]{\textcolor{red}{[TODO: #1]}}

\graphicspath{{figures/}}

\journal{Robotics and Autonomous Systems}

\begin{document}

\begin{frontmatter}

\title{Resilient Collaborative Perception for Multi-Robot Navigation under
Sensor Noise and Byzantine False-Obstacle Attacks}

%% Authors — fill affiliations/emails before submission
\author[inst1]{K.~Srinivasa~Rao\corref{cor1}}
\ead{raokommanasrinivasa@gmail.com}
\author[inst1]{\todo{Professor name}}
\cortext[cor1]{Corresponding author.}
\affiliation[inst1]{organization={\todo{Institute name, department}},
            city={\todo{city}},
            country={India}}

\begin{abstract}
% <= 250 words, standalone, no references, no undefined abbreviations.
% DRAFT v0 — update from PAPER_MASTER_PLAN sec 9.1, with CAMERA-READY numbers
% (500 maps, density 0.27, recall 0.13->0.69, recovery +3.4->+12.2 pp).
Multi-robot teams can survive individual sensor failure by sharing perception
over a communication channel, but that same channel is an attack surface: a
single Byzantine teammate broadcasting fabricated obstacles can override
honest on-board sensing through minimum-based map fusion, independently of
the sensor-failure rate. We study this threat in a 10-drone navigation task
learned with centralized-training/decentralized-execution multi-agent
reinforcement learning, and show that the natural defense --- a cross-agent
consistency check --- is brittle: under realistic ranging noise it wrongly
condemns honest neighbours and becomes \emph{worse than no defense}, and even
a noise-aware robust variant fails against a camouflage attack that hides the
fabrication inside the sensor-noise band (detection recall 0.13). We
introduce a temporal trust rule that aggregates, over time, the offset
between a neighbour's reported obstacle and the verifier's own view of it:
honest disagreement is zero-mean and cancels, while a persistent fabrication
retains a bias. This recovers detection recall to 0.69 and navigation success
by up to +12.2 percentage points at the severest noise level, at no
measurable cost to attack-free swarms. Against an adaptive, defense-aware
attacker we demonstrate a stealth/harm bind: evading temporal detection
forces the phantom onto a real obstacle, where it blocks no new space.
Because each verdict is pairwise rather than a majority vote, the mechanism
holds where consensus-based defenses cannot: the recovery remains
statistically significant with up to seven of ten robots lying, past the
point where an honest robot's neighbourhood is majority traitor.
Results use 500-map evaluation suites with paired-bootstrap confidence
intervals. % ~240 words

\end{abstract}

%% 1-7 keywords, single words preferred
\begin{keyword}
multi-robot systems \sep collaborative perception \sep Byzantine resilience
\sep trust \sep reinforcement learning \sep swarm navigation
\end{keyword}

\end{frontmatter}

% ---------------------------------------------------------------------
\section{Introduction}\label{sec:intro}
% Sources: PAPER_MASTER_PLAN sec 9 (structure), RESULTS_027_CAMERA_READY.md
% (all numbers). Camera-ready figures: 500 maps, density 0.27, stage-2 base.

Teams of small autonomous robots are increasingly deployed for tasks in
which a single vehicle is insufficient: warehouse logistics, search and
rescue, environmental monitoring, and infrastructure inspection. In
cluttered environments each robot perceives only a fraction of its
surroundings --- on-board range sensors are limited in reach and are
routinely degraded by occlusion, reflective or absorptive surfaces, dust,
and transient hardware faults. \emph{Collaborative perception} addresses
this by letting teammates share what they sense, so that one robot's view
fills another's blind spot. In the navigation benchmark studied here --- ten
drones crossing a cluttered field while suffering sustained sensing
blackouts roughly one third of the time --- sharing is not a refinement but
a necessity: with collaborative perception 89.3\% of drones reach the goal,
without it only 45.9\% do, a gap of more than 43 percentage points under
identical conditions (Section~\ref{sec:results}).

The channel that carries this resilience, however, is itself an attack
surface. Shared observations must be fused into each robot's local obstacle
map, and safe fusion is conservative: if any teammate reports an obstacle
nearby, the robot should respect it. We show that this conservatism ---
implemented here by a per-direction \emph{minimum} over own and
received range estimates, the choice that respects any teammate's evidence
of a nearby obstacle --- hands a powerful lever to a compromised
insider. A single \emph{Byzantine} teammate that broadcasts fabricated
obstacles overrides even a fully-sighted receiver's own sensing, because a
fabricated near reading always wins the minimum. The attack is therefore
independent of the sensor-failure rate that motivates collaboration in the
first place: it succeeds against healthy sensors as easily as against
degraded ones. Two traitors among ten drones cut honest navigation success
by up to 22.5 percentage points, and a \emph{camouflage} variant that
extends a real obstacle into a free corridor is both more damaging and far
harder to detect than a blatant phantom wall.

The natural defense is cross-agent consistency checking: if a verifier is
positioned to see a reported obstacle and does not, the reporter loses
trust and is excluded from fusion. Under idealized sensing this neutralizes
the attack almost completely. It breaks down, however, in exactly the
conditions that make collaboration valuable. Under realistic ranging noise
two honest views of the same obstacle legitimately disagree by about
$\sqrt{2}\sigma$; once that exceeds the verifier's fixed tolerance, honest
neighbours are condemned as liars. Detection precision collapses from
$1.00$ to $0.23$, and merely enabling the filter destroys up to 27
percentage points of success in swarms containing no attacker at all:
\emph{the defense is worse than no defense}. Widening the tolerance with
$\sigma$ repairs the false accusations but opens the complementary hole ---
a camouflage attacker hides its fabrication \emph{inside} the widened band,
where no single-frame test can contradict it, and at the severest noise
level such a filter catches only 13\% of camouflage lies. The bind is
structural rather than a matter of tuning: on any single frame, a tolerance
tight enough to expose the lie is tight enough to condemn an honest
neighbour. That this difficulty is not hypothetical is borne out by an early
study of trust-modulated cooperative perception, which found that
cooperative fusion did not meaningfully improve accuracy over local
perception, with false detections inserted at high confidence among the
suspected causes~\cite{trupercept2020}.

Our central contribution closes this blind spot with time rather than with
looser thresholds. The per-frame \emph{offset vector} between a neighbour's
reported obstacle and the verifier's own view of that obstacle is zero-mean
for an honest neighbour --- two independent noise draws cancel in
expectation --- but retains a \emph{persistent bias} for a fabrication,
whose reported position is anchored to a lie rather than to the shared
ground truth. Aggregating this offset over a few tens of observations
separates the two populations cleanly. This use of time differs
fundamentally from recent temporal-consistency defenses for cooperative
perception, which compare a scene against its own recent past and so screen
for \emph{disruption}: once established, a persistent fabrication can
remain temporally self-consistent --- a signature such defenses do not
study --- whereas it leaves a persistent offset against the verifier's
\emph{own sensing}, which is exposed through accumulation over time
(Section~\ref{sec:related}). The resulting temporal trust rule
lifts camouflage detection recall from $0.13$ to $0.69$ and recovers up to
$+12.2$ percentage points of navigation success at the severest noise
level, while leaving attack-free swarms statistically untouched (the
measured cost is $-0.4$ percentage points with a confidence interval
spanning zero). Because each verdict is pairwise --- a robot tests each
neighbour against its own sensing and never counts votes --- the rule needs
no honest majority: its recovery remains statistically significant with up
to seven of ten robots lying, past the point where an honest robot's
neighbourhood is majority traitor, whereas the single-frame check is no
longer distinguishable from no defense there. Anticipating the strongest objection --- that a static
attacker is a strawman --- we also evaluate an adaptive attacker that knows
the defense and tunes its fabrication to evade it. The attacker faces a
\emph{stealth/harm bind}: as the phantom drifts far enough from real
geometry to block useful space, its offset signature grows and detection
recall rises with the harm; hiding from the filter forces the phantom onto
a real obstacle, where it blocks nothing new. We verify this bind at every
noise level and at one to three traitors.

This paper makes five contributions. \textbf{(i)}~We quantify the
resilience value of collaborative perception for learned multi-robot
navigation under sustained sensor dropout at a calibrated obstacle density,
with paired-bootstrap confidence intervals over 500-map suites.
\textbf{(ii)}~We characterize a minimum-fusion Byzantine vulnerability of
shared obstacle maps and show it is independent of the dropout rate.
\textbf{(iii)}~We show that naive consistency-based trust is not merely
weakened but \emph{destructive} under sensor noise. \textbf{(iv)}~We give a
noise-aware robust filter that restores safe operation and characterize its
single-frame camouflage limit. \textbf{(v)}~We introduce a temporal
offset-bias trust rule that recovers this limit at no measured cost to
honest swarms and, being pairwise, retains a significant recovery with up to
seven of ten robots lying --- past the loss of an honest local majority ---
and we validate it against an adaptive, defense-aware attacker
across one to three traitors, demonstrating a stealth/harm bind that makes
the recovery difficult to game.

The remainder of the paper is organized as follows.
Section~\ref{sec:related} positions the work against collaborative
perception security, temporal spoof detection, and Byzantine multi-robot
systems. Section~\ref{sec:methods} details the environment, the learned
navigation policy, the threat model, and the three defenses.
Section~\ref{sec:setup} defines metrics and the statistical protocol.
Section~\ref{sec:results} reports results, Section~\ref{sec:discussion}
discusses limitations honestly, and Section~\ref{sec:conclusion} concludes.


\section{Related work}\label{sec:related}
% Sources: PAPER_MASTER_PLAN sec 9.2 (audit 2026-06-19) + sec 10.1 re-sweep
% (2026-06-26). Verify TODO-VERIFY bib entries before submission.

\paragraph{Security of collaborative perception.}
Attacks on, and defenses of, shared perception have been studied intensively
for connected autonomous vehicles. Zhang et al.~\cite{zhang2024cad}
demonstrate practical data-fabrication attacks on vehicular collaborative
perception and propose CAD, a detector in which benign vehicles jointly
reveal malicious fabrication through cross-vehicle occupancy-consistency
checks; the
approach therefore relies on benign vehicles positioned to observe the
attacked region. Two aspects of our setting go beyond this. First, we do not
assume such a benign observer of the attacked region: our test is pairwise
and remains effective with up to three traitors among ten robots. Second, we
specifically target \emph{camouflaged} fabrications --- phantoms blended into
real geometry --- which are the hardest to separate from honest observation
error and which a single-frame cross-agent check cannot contradict under
sensor noise. In feature-fusion collaborative perception,
GCP~\cite{gcp2025} defends against malicious agents by augmenting
single-shot spatial-consistency checks with temporal motion-flow
reconstruction to expose blind-area attacks; MADE~\cite{made2024} detects
and removes malicious agents through hypothesis tests on the consistency
between each inspected agent and the ego agent, with statistical
false-positive control; and CoDynTrust~\cite{codyntrust2025} addresses the
benign counterpart by weighting intermediate feature contributions
according to their estimated uncertainty, thereby mitigating temporal
asynchrony rather than adversarial behaviour; and ROBOSAC~\cite{robosac2023}
rejects adversarial feature-map perturbations by sampling subsets of teammates
until the collaborative output reaches consensus with the ego's own perception,
falling back to ego-only perception when no consensus is found. All four are
evaluated using Average Precision (AP) on vehicular benchmarks.

A second, closer line verifies \emph{object-level claims against the
verifier's own observations}. TruPercept~\cite{trupercept2020} evaluates reported
object detections as verified locally, weighting each evaluation by the
visibility of the claimed object from the verifier's viewpoint before
computing peer trust; MATE~\cite{mate2025} formalizes this as Bayesian trust
estimation over geometric observation consistency with explicit
field-of-view reasoning, against false-positive, false-negative, and
translation attacks on multi-target tracking; and the same authors' aerial
framework~\cite{aerialtrust2025} carries that hidden-Markov trust estimation
to ad hoc UAV networks on intelligence, surveillance, and reconnaissance
(ISR) missions --- the setting nearest to ours in the literature. Our single-frame filters (Section~\ref{sec:methods-defense})
belong to precisely this family, and inherit its central assumption: that a
verifier positioned to observe a claimed object either corroborates or
contradicts it.

That assumption is what a camouflaged fabrication attacks. It is instructive
that TruPercept additionally rejects \emph{open-space} phantoms with a
plausibility check that culls the sensor frustum toward a claimed object: an
object with no returns in front of it, but returns behind it, is deemed
absent with high confidence~\cite{trupercept2020}. Operationally, such a
test derives its evidence from rays passing through empty space at the
claimed location, which is exactly what a phantom placed in free space provides ---
and, in our own experiments, one reason the conspicuous wall attack proves
comparatively easy to detect (Section~\ref{sec:res-temporal}). A camouflaged
phantom is designed to avoid supplying that evidence: placed flush against a
real obstacle, its claimed volume coincides with returns the verifier
genuinely receives from nearby structure, so a free-space plausibility test
has little to discriminate on. We note that TruPercept does not evaluate
camouflaged fabrications, and we make no claim about the behaviour of its
checker in our setting; the point is that geometric plausibility tests
derive their power from rays traversing empty space, and camouflage attacks
intentionally deny them that. Under ranging noise the residual discrepancy
further shrinks below the tolerance a verifier must grant honest neighbours
(Section~\ref{sec:res-robust}). Our temporal test therefore does not seek a
stronger \emph{spatial} contradiction; it accumulates a \emph{statistical}
one over time.

This line accommodates natural sensor error through longitudinal filtering,
and observes that an attacker must remain dynamics-consistent to evade local
filtering; TruPercept further reports that unreliable agents which mostly
report truthfully are not separated from honest ones by its trust
model~\cite{trupercept2020} --- an early indication of the evasion regime we
formalize. Where our setting departs further:
TruPercept~\cite{trupercept2020} aggregates trust at a central server and
MATE~\cite{mate2025} at a central computing centre (the aerial
framework~\cite{aerialtrust2025} does distribute the estimation), whereas
our verdicts are pairwise and local, as a swarm without infrastructure
requires. None of this line
characterizes the regime in which the consistency check \emph{itself} turns
harmful under ranging noise (our destructive-filter result), considers
fabrications geometrically camouflaged within the noise tolerance of real
objects, or measures impact on a closed-loop learned \emph{navigation} task
--- their trust protects a shared perception picture, whether for driving, a
smart city, or aerial surveillance, whereas ours protects the obstacle map a
policy steers by. Our contribution begins where this family breaks. Closest to our mechanism,
the recent PRBI defense~\cite{prbi2026} exploits frame-to-frame perceptual
consistency to identify lying vehicles in fully untrusted cooperative
detection. It differs from our setting in three ways that matter: it is
evaluated by detection accuracy on feature-level fusion rather than by a
closed-loop navigation objective; it does not model ranging noise, and so
does not confront the honest-disagreement regime in which we show
consistency checking becomes destructive; and while it evaluates
intermittent injection schedules, it does not consider a defense-aware
attacker that optimizes its fabrication against the deployed detector, as
our stealth/harm-bind analysis does. More fundamentally, PRBI's alarm is
raised when the Jaccard similarity between the current and preceding
detection sets falls below a threshold~\cite{prbi2026}: it screens for a
\emph{drop} in inter-frame agreement. Such a criterion responds to
perturbations of the scene relative to its recent past. A
persistent fabrication is temporally self-consistent once established;
PRBI does not evaluate this attack class, and we make no claim about how
its detector would behave against it. The
distinction we draw is one of reference signal: PRBI compares the scene to
its own past, whereas our test compares a neighbour's claim to the
verifier's independent sensing of the same object, accumulated over time.

Complementing the defense literature, the TrustFlip
attack~\cite{trustflip2026} shows that consistency-based trust can itself be
weaponized: physically induced disagreements cause defenses to expel
\emph{honest} vehicles. Our evaluation addresses the same failure class
by measuring \emph{no-harm} --- the success cost of running the defense in
attack-free swarms --- directly, and finding it statistically
indistinguishable from zero even under an adaptive attacker.

\paragraph{Temporal detection of sensor spoofing.}
A separate line detects LiDAR spoofing against a \emph{single} vehicle by
exploiting temporal structure in its own sensor stream, e.g.\ motion-induced
consistency in 3D-TC2~\cite{tc2_2021} and point-level temporal consistency
in ADoPT~\cite{adopt2023}. These methods test whether observations from a
single sensor remain temporally self-consistent. Ours instead aggregates a \emph{cross-agent
disagreement} --- the offset between a neighbour's claim and the verifier's
own view --- whose statistics under honest noise (zero-mean, variance
$2\sigma^2$) versus fabrication (persistent bias) we can characterize and
threshold explicitly. The statistic itself is elementary; the contribution
is its use as a communication-trust discriminator that survives sensor
noise and camouflage, conditions under which we show single-frame
cross-agent checks fail.

\paragraph{Byzantine resilience in multi-robot systems.}
Byzantine fault tolerance originates in distributed
computing~\cite{lamport1982byzantine} and, in robot swarms, is typically
addressed at the \emph{consensus} layer. Consensus-based coordination such
as SwarmRaft~\cite{swarmraft2025} fuses peer measurements to maintain
agreement on state (e.g.\ position and heading) under degraded conditions,
and evolutionary-game analyses~\cite{conformity2026} model how deceptive
strategies propagate through conformity dynamics under Byzantine influence.
These works protect state agreement or model attacks on collective
decision-making; our
threat lives one layer below, in the \emph{perception content} that robots
exchange, and our defense verifies that content against the verifier's own
sensing rather than against majority agreement. Because each verdict is
pairwise, the mechanism is structurally independent of the number of liars;
we confirm this empirically (Section~\ref{sec:res-temporal}), where its
recovery remains significant even when traitors form a local majority among
an honest robot's neighbours --- up to seven of ten robots. Sampling-based
consensus defenses such as ROBOSAC~\cite{robosac2023}, by contrast, must draw
an attacker-free subset of teammates and know or estimate the attacker ratio,
with a sampling budget that grows sharply as benign agents become scarce; our
pairwise test needs neither a clean subset, a known ratio, nor a consensus
vote.

\paragraph{Learned multi-robot navigation and robust communication.}
Our navigator is trained with centralized-training/decentralized-execution
multi-agent reinforcement learning, following the MAPPO
recipe~\cite{yu2022mappo} built on PPO~\cite{schulman2017ppo} and, in
implementation, on Stable-Baselines3~\cite{raffin2021sb3}; CTDE itself
traces to multi-agent actor-critic formulations such as
MADDPG~\cite{lowe2017maddpg}. Our attack surface is deliberately
interpretable: the messages are geometric obstacle claims rather than
learned latent vectors, the fusion rule is explicit, and the defense is a
hand-designed statistical test --- one of our findings is that no learned
trust component is required to counter this threat class.

The closed-loop learned cooperative-perception paradigm we build on --- a
policy that consumes shared perception and is scored by end-task success ---
was established for benign networked driving by
Coopernaut~\cite{coopernaut2022}; our contribution is not that paradigm but
the adversarial layer within it. Closest in driving-level outcome, the recent
PosePert attack~\cite{stealthyfab2026} propagates small, stealthy pose
perturbations of \emph{existing} objects through a modular tracking-and-
prediction stack to induce unsafe behaviour on a single vehicle, measured by
trajectory-prediction error and a scripted braking-risk rate. Our setting
differs in the attack (fabricated obstacles, not pose shifts of real objects),
the control layer (a learned navigation policy executing in closed loop, rather
than a tracker-plus-predictor pipeline), the end metric (task success, not
prediction error or a safety flag), and the multi-robot rather than
single-vehicle scope.

In combination, to our knowledge no prior work evaluates
fabricated-obstacle attacks on collaborative perception \emph{inside a
closed-loop learned navigation task} (success, not detection accuracy, as
the end metric), analyzes the interaction of consistency-based trust with
sensor noise to the point of demonstrating harmfulness, closes the
resulting camouflage blind spot with a cross-agent temporal bias test, and
stress-tests the closure against an adaptive attacker across noise levels
and traitor counts.

\paragraph{On baseline comparison.}
The defenses above cannot be transplanted onto our benchmark without
distortion: CAD, PRBI, GCP, and MADE defend the fused perception pipelines
of vehicular detection stacks and are evaluated by detection-level metrics
(anomaly-detection rates or average precision), while TruPercept, MATE, and
the same authors' aerial framework operate at the object and track level;
none is designed to operate directly on the geometric obstacle claims
exchanged by our agents, and none produces the navigation-success outcome
we measure --- a reimplementation would test our translation of each method
rather than the method itself. We
instead compare against the \emph{primitive that the object-level family
shares}: verifying a peer's claims against what the verifier itself was
positioned to observe, accumulating trust, and excluding on contradiction.
That primitive is implemented natively in our setting in two strengths ---
the fixed-tolerance (naive) variant and the noise-aware robust variant
(Section~\ref{sec:methods-defense}) --- and it is the same
visibility-and-consistency recipe underlying
TruPercept~\cite{trupercept2020} and the MATE
line~\cite{mate2025,aerialtrust2025}. Our evaluation therefore reports how
the family's core mechanism behaves in this regime --- including its
noise-induced failure --- before showing what the temporal test adds on top
of it.


\section{System model and methods}\label{sec:methods}
% Verified against code 2026-07-08:
%   env_noisy_byzantine.py (_sample_sensing l.144, _ego_judgement l.170,
%   _temporal_update l.110, _broadcast_phantoms l.87, _fused_lidar l.186).
% Numbers/params: CLAUDE.md, PARAMETER_JUSTIFICATION_PHASE_CD.md, M0_PROVENANCE.
% NOTE: verify_eps=0.6, k_sigma=4, alpha=0.25, tau=0.4, min_k=20, bias_eps=0.6.

\subsection{Task and environment}\label{sec:methods-env}
We consider $N=10$ holonomic drones navigating a shared $20\times20$~m
planar field to a common goal. The field is populated with static circular
obstacles at a calibrated density (Section~\ref{sec:setup}) chosen so that
the task is neither trivial nor infeasible; every evaluated map is verified
solvable. Each drone observes its surroundings through a 48-ray planar
LiDAR of range $R=8$~m and selects a continuous velocity action; episodes
run for up to 1200 steps. Drones that reach the goal are removed from the
scene so that they neither obstruct nor collide with the remainder.

Perception is deliberately \emph{partial}: with an $8$~m range on a $20$~m
field, no drone observes the whole environment, which is the precondition
for collaboration to be meaningful. Drones may exchange messages within a
communication range of $10$~m, strictly larger than the sensing range so
that a neighbour's message can describe regions the receiver cannot yet
sense itself; were communication no farther than sensing, sharing could
only ever duplicate what the receiver already sees. An external mission
planner supplies each drone with a routed goal-direction heading, computed
as the gradient of a global shortest-path field; the learned policy performs
local control and collision avoidance given this heading. We treat the
routed heading as a disclosed input rather than a learned capability
(Section~\ref{sec:discussion}).

\paragraph{Scope of the abstraction.}
We deliberately model perception at the \emph{object level}: agents
exchange geometric obstacle claims (centre, radius) rather than raw point
clouds or learned feature maps. This abstraction is chosen, not incidental.
The questions this paper asks --- when does a fabricated claim override
honest sensing, when does consistency checking turn harmful under noise,
and what evidence separates a persistent lie from honest error --- concern
the \emph{fusion and trust layer}, and posing them at the object level
isolates that layer from the particulars of any perception stack, in the
same way that classical multi-target-tracking analyses reason over
detections rather than pixels. The price is that phenomena living below
this level --- occlusion patterns, detector biases, feature-space attacks
--- are out of scope here (Section~\ref{sec:discussion}); the benefit is
that the trust mechanisms and their failure modes are characterized
independently of a specific sensor pipeline, and the offset-bias statistic
we exploit is defined for any system that can associate a neighbour's
claimed object with its own view of that object.

\subsection{Learned navigation policy}\label{sec:methods-policy}
The navigation policy is trained with multi-agent proximal policy
optimization under the centralized-training/decentralized-execution
paradigm~\cite{yu2022mappo,schulman2017ppo}. During training a centralized
critic observes a $520$-dimensional global state, while each actor observes
only a $130$-dimensional local vector --- its own LiDAR, kinematic state,
routed heading, and gated neighbour summaries. At execution only the
decentralized actor runs; the critic is discarded. The policy is obtained by
curriculum transfer over increasing obstacle density and, for the
experiments reported here, a final fine-tuning stage under sensor-noise
domain randomization ($\sigma\sim\mathcal{U}[0,0.6]$~m) so that the base
navigator is not disadvantaged at test time by noise it never saw in
training. Full hyperparameters and the training lineage are provided in the
supplementary material; they are not central to the defense contribution.

\subsection{Collaborative perception and fusion}\label{sec:methods-fusion}
Each sighted drone $j$ senses the set of obstacles within its range and
broadcasts their estimated centres and radii to neighbours in
communication range. A receiver $i$ fuses received obstacles with its own by
rendering all of them into its 48-ray LiDAR channel and taking, per ray,
the \emph{minimum} distance over own and received returns. Minimum-fusion
is one point in a design space: averaging or confidence-weighted
combination would dilute both fabricated and genuine individual returns,
potentially weakening the property that a single teammate's warning is
sufficient to trigger avoidance. We adopt the conservative rule and defend
it, rather than trading safety for attack resistance; the vulnerability
studied below is the price of that safety property, not an artifact of an
eccentric choice.

We model such blackouts explicitly. At each step a drone independently
enters a sustained blind burst with probability $p=0.10$ lasting $s=5$
steps, giving a stationary blind fraction of $ps/(1+ps)\approx 0.33$; a
blind drone senses nothing and, crucially, \emph{broadcasts nothing}
(sender gating), so comm cannot trivially substitute for a drone's own
missing sensor. This regime is where sharing is load-bearing: it is the
operating point at which the gap between sharing and not sharing is largest
while the map remains solvable (Section~\ref{sec:results}).

\subsection{Threat model: Byzantine false-obstacle attacks}\label{sec:methods-threat}
Among the $N$ drones, $f\in\{1,2,3\}$ are Byzantine insiders. They obey the
communication protocol but broadcast \emph{fabricated} obstacles that do not
exist in the ground truth. Because fabricated obstacles are not physically
present, the map remains solvable --- the attack denies passage by deception
rather than by blocking it, so any lost success is attributable to the lie,
not to an impossible task. A fabricated near obstacle always wins the
per-ray minimum in fusion, so a single traitor overrides even a
fully-sighted receiver's own correct sensing; the attack is consequently
\emph{independent of the dropout rate} that motivates collaboration.

We study two placements. In the \emph{wall} attack the traitors broadcast a
line of phantom obstacles across the goal approach --- damaging but
conspicuous, since a sighted honest drone is positioned to notice that the
claimed barrier is not there. In the \emph{camouflage} attack each phantom
is placed hugging a real obstacle bordering a free corridor, extending it
into the gap; this is both more harmful (it closes the route a drone would
actually take) and far stealthier (the phantom sits where obstacles
plausibly are). To avoid tuning the attack to a single easy configuration,
the fabrication is \emph{randomized} per map: the number of phantoms is
drawn as $n\sim\mathcal{U}\{3,4,5,6\}$ and each phantom radius is drawn
from a three-class sampler ($40/40/20\%$ small/medium/large over the same
radius ranges as the map generator), matched to the measured size
distribution of the real obstacles ($42/40/18\%$), so phantoms are
size-indistinguishable from genuine obstacles and cannot be filtered on
size alone.

We deliberately grant the attacker more than is physically realizable: it
places phantoms using the \emph{true} obstacle map and broadcasts their
positions noise-free, while honest drones sense and report under noise.
Reported damage is therefore measured against a best-case adversary. The
adaptive attacker of Section~\ref{sec:res-adaptive} is
\emph{design-aware but not state-aware}: it knows how the deployed defense
works and tunes its fabrication geometry against it, but it cannot observe
its own runtime trust scores or other drones' private observations, and so
cannot adapt within an episode.

\subsection{Sensor-noise model}\label{sec:methods-noise}
Real range sensing is imperfect. We add independent zero-mean Gaussian noise
of standard deviation $\sigma$ to every drone's estimate of every sensed
obstacle position, resampled each step:
$\hat{c}_{j,m}(t) = c_m + \varepsilon_{j,m}(t)$, with
$\varepsilon_{j,m}(t)\sim\mathcal{N}(0,\sigma^2 I_2)$, independent across
drones $j$, obstacles $m$, and time $t$. We sweep
$\sigma\in\{0,0.2,0.4,0.6\}$~m. At $\sigma=0.6$ the positional error is on
the order of the obstacle radius itself (mean $0.91$~m), several times the
drone clearance, and --- the key consequence for defense --- large enough
that two honest drones sensing the same obstacle disagree by
$\lVert\varepsilon_{i,m}-\varepsilon_{j,m}\rVert$ with expected magnitude on
the order of $\sqrt{2}\,\sigma\approx0.85$~m.

\subsection{Defenses}\label{sec:methods-defense}
Every defense maintains a per-pair trust scalar $t_{ij}\in[0,1]$, reset to
$1$ each episode. When receiver $i$ forms a judgement about neighbour $j$ on
a given step it updates trust by exponential moving average,
\begin{equation}
t_{ij}\leftarrow(1-\alpha)\,t_{ij}+\alpha\,\mathbb{1}[\text{$j$ not
contradicted}],\qquad \alpha=0.25,
\label{eq:ewma}
\end{equation}
and excludes $j$ from fusion whenever $t_{ij}<\tau$ ($\tau=0.4$). With
$\alpha=0.25$, trust builds gradually under consistent behaviour but
collapses within a few contradictions; the same qualitative requirement ---
gradual accrual, rapid degradation --- has been articulated independently,
within a different (Bayesian) trust formalism, by Hallyburton and
Pajic~\cite{aerialtrust2025}. The three
defenses differ only in how a contradiction is decided.

\paragraph{(a) Naive single-frame trust.}
If $i$ is sighted and positioned to see a broadcast obstacle $o$ (i.e.\
$\lVert p_i-o\rVert\le R$) but its own nearest sensed obstacle lies farther
than a fixed tolerance $\varepsilon_0=0.6$~m from $o$, then $j$ is
contradicted on that frame. Under noiseless sensing this excludes traitors
almost perfectly. Under noise it is fragile: two honest views of one
obstacle differ by $\sim\sqrt{2}\sigma$, which for $\sigma\ge0.4$~m exceeds
$\varepsilon_0$, so honest neighbours are contradicted and gated out
(Section~\ref{sec:results}).

\paragraph{(b) Noise-aware robust trust.}
The robust filter widens the tolerance in proportion to the known noise
level,
\begin{equation}
\varepsilon(\sigma)=\varepsilon_0+k_\sigma\,\sigma,\qquad k_\sigma=4,
\label{eq:robusteps}
\end{equation}
which restores $\approx2\sqrt{2}\sigma$ coverage of honest disagreement with
margin, so genuine noisy matches pass while phantoms in open space (metres
from any real obstacle) still fail. This repairs the false-accusation
problem, but a camouflage phantom placed within $\varepsilon(\sigma)$ of the
real obstacle it hugs cannot be contradicted on any single frame: the lie
hides inside the widened band.

\paragraph{(c) Temporal offset-bias trust.}
Our contribution attacks the camouflage phantom over time rather than with a
looser threshold. For each broadcast obstacle $o$, receiver $i$ associates
it with its own nearest sensed obstacle --- call the matched track $m$ and
$i$'s current noisy view of it $\hat{c}_{i,m}$ --- and records the
\emph{offset vector}
\begin{equation}
d = o - \hat{c}_{i,m}.
\label{eq:offset}
\end{equation}
For an honest neighbour reporting the true obstacle,
$o=\hat{c}_{j,m}=c_m+\varepsilon_{j,m}$, so
$d=\varepsilon_{j,m}-\varepsilon_{i,m}$ is zero-mean with variance
$2\sigma^2$: averaged over frames it \emph{cancels}. For a fabrication,
$o$ is anchored to the lie rather than to $c_m$, so $d$ carries a
\emph{persistent bias} equal to the phantom's displacement from the matched
real obstacle, which does \emph{not} cancel. The receiver keeps a running
mean of $d$ in a bucket indexed by $(i,j,m)$; once a bucket has accumulated
at least $K_{\min}=20$ samples it yields a verdict, contradicting $j$ if the
mean offset magnitude exceeds a threshold,
\begin{equation}
\Big\lVert\tfrac{1}{n}\textstyle\sum_{t} d_t\Big\rVert > \eta,
\qquad \eta=0.6~\text{m},\; n\ge K_{\min}.
\label{eq:temporal}
\end{equation}
Because $\eta$ is applied to a \emph{mean} rather than a single frame, it can
stay tight while honest noise (which shrinks as $1/\sqrt{n}$) falls below it.
The temporal test composes with the single-frame robust test by logical OR:
$j$ is contradicted if either fires, so the fast path catches conspicuous
wall phantoms immediately while the slow path catches camouflage phantoms
once enough evidence accrues. Algorithm~\ref{alg:temporal} summarizes one
step of the composed defense. All thresholds
($\varepsilon_0,k_\sigma,\alpha,\tau,K_{\min},\eta$) are justified
individually in Section~\ref{sec:setup} and the supplementary material; none
is tuned per condition. Two implementation notes: the receiver's identity of
its \emph{own} tracks $m$ is idealized (obstacles are static, so persistent
track identity is mild), whereas the broadcast-to-track association in
Eq.~\eqref{eq:offset} is realistic nearest-neighbour matching --- a probe of
the offset signal alone yields ROC-AUC $0.99$ under oracle association and
$0.85$--$0.88$ under this realistic association (both over 500 maps),
quantifying the cost of the weaker assumption. The filter is lightweight: the running means occupy under
$100$~KB for the whole swarm, and the measured overhead of the composed
defense is ${\sim}0.15$~ms per drone per step in unoptimized Python.

\begin{algorithm}[t]
\DontPrintSemicolon
\caption{Composed trust defense, one step at receiver $i$.}
\label{alg:temporal}
\KwIn{own noisy obstacle views $\hat{c}_{i,\cdot}$; for each in-range
neighbour $j$: broadcast obstacles $O_j$; trust $t_{ij}$; noise $\sigma$}
\KwParam{$\varepsilon_0{=}0.6$, $k_\sigma{=}4$, $\alpha{=}0.25$,
$\tau{=}0.4$, $K_{\min}{=}20$, $\eta{=}0.6$}
$\varepsilon \leftarrow \varepsilon_0 + k_\sigma\,\sigma$\;
\ForEach{in-range neighbour $j$}{
  $\mathit{judged}\leftarrow\mathbf{false}$;\quad
  $\mathit{bad}\leftarrow\mathbf{false}$\;
  \ForEach{$o \in O_j$ with $\lVert p_i-o\rVert \le R$}{
    $m \leftarrow \arg\min_{m'} \lVert \hat{c}_{i,m'} - o\rVert$\quad
      \tcp*{match to own nearest obstacle}
    $\mathit{judged}\leftarrow\mathbf{true}$\;
    \lIf{$\lVert \hat{c}_{i,m} - o\rVert > \varepsilon$}{
      $\mathit{bad}\leftarrow\mathbf{true}$ \tcp*[f]{single-frame robust test}}
    $S_{ijm}\mathrel{+}= (o-\hat{c}_{i,m})$;\quad $n_{ijm}\mathrel{+}=1$\;
    \If{$n_{ijm}\ge K_{\min}$ \textbf{and}
        $\lVert S_{ijm}/n_{ijm}\rVert > \eta$}{
      $\mathit{judged},\mathit{bad}\leftarrow\mathbf{true},\mathbf{true}$
        \tcp*{temporal offset-bias test}
    }
  }
  \If{$\mathit{judged}$}{
    $t_{ij}\leftarrow (1-\alpha)\,t_{ij}
      + \alpha\,(0 \text{ if } \mathit{bad} \text{ else } 1)$\;
  }
  \lIf{$t_{ij} < \tau$}{exclude $O_j$ from fusion}
}
\end{algorithm}


\section{Experimental setup}\label{sec:setup}
% Sources: RESULTS_027_CAMERA_READY.md (setup block), boot_ci.py,
% measure_env_stats.py, PARAMETER_JUSTIFICATION_PHASE_CD.md. Verified.

\subsection{Metrics}\label{sec:setup-metrics}
Our primary metric is \emph{honest-drone success}: the fraction of
non-Byzantine drones that reach the goal within the episode horizon,
averaged over a map suite. We report it as a percentage. We also track
\emph{map-level success} (all ten drones reach the goal) for the
collaborative-perception anchor, where it is the more demanding and more
intuitive figure. For the defenses we additionally report, for each
condition:
\begin{itemize}
\item \textbf{Recovery} $=$ (success with the defense on, under attack) $-$
(success with no defense, under attack), in percentage points; positive
means the defense helps.
\item \textbf{No-harm} $=$ (success with the defense on, \emph{no} attacker)
$-$ (attack-free base success). This isolates the cost of running the
defense itself; a value near zero means the defense does not wrongly gate
honest neighbours. It is the direct measure of the failure mode that
consistency defenses are known to suffer, and we treat it as a
first-class result rather than a proxy.
\item \textbf{Detection precision and recall} of the trust test, where a
true positive is a correctly contradicted traitor link and a false positive
is a wrongly contradicted honest link.
\end{itemize}

\subsection{Evaluation protocol}\label{sec:setup-protocol}
Each condition is evaluated on \textbf{500 maps}. Obstacles are circular,
placed at a target density of $0.27$; at this density a
solvability-verified map contains on average $29.7$ obstacles (range
$15$--$56$) with a probability-weighted mean radius of $0.91$~m, and
$96.8\%$ of sampled maps are solvable. Maps are generated from
deterministic seeds so that every arm of a comparison --- base, attack,
robust, temporal, no-harm --- is evaluated on the \emph{identical} suite of
maps. This paired design lets us report \textbf{paired-bootstrap $95\%$
confidence intervals}: we resample map indices with replacement (2000
resamples, fixed seed) and recompute each quantity, taking the $2.5$th and
$97.5$th percentiles; for a difference such as recovery or no-harm we
resample the per-map \emph{difference}, which cancels shared map-difficulty
variance and yields tight intervals. An interval that excludes zero
indicates a statistically reliable effect.

\subsection{Conditions}\label{sec:setup-conditions}
The main defense study sweeps noise $\sigma\in\{0,0.2,0.4,0.6\}$~m, attack
placement $\in\{\text{wall},\text{camouflage}\}$, and traitor count
$f\in\{1,2,3\}$ (a bare minority, the literature-standard ceiling), each
under five arms (base, attack, robust, temporal, no-harm). The adaptive
attacker is studied with four independent evasion knobs --- phantom
centre-offset, surface gap, per-frame jitter, and broadcast duty --- swept
at $\sigma=0.6$; because these sweeps isolate a single attacker capability,
they hold the phantom radius fixed rather than randomizing it, so the swept
axis is not confounded. The offset knob, which cleanly exhibits the
stealth/harm bind, is additionally swept across all three noise levels and
all three traitor counts.

\subsection{Reproducibility}\label{sec:setup-repro}
The simulation environment, the trained navigation and collaborative-
perception models, all evaluation scripts, and the raw per-condition result
logs (including the confidence-interval outputs) are released publicly so
that every number in Section~\ref{sec:results} can be regenerated with a
single driver script. The release and its archival identifier are given in
the data-availability statement.


\section{Results}\label{sec:results}
% ALL NUMBERS from RESULTS_027_CAMERA_READY.md (camera-ready: 500 maps,
% density 0.27, stage-2 base, randomized attack, paired-bootstrap 95% CIs).
% Probe AUC now 500-map (0.99 oracle / 0.85-0.88 realistic, 0.88 at K_min=20)
% -- rerun done 2026-07-20, rerun/probe_{oracle,realistic}_500.txt; confirms
% the earlier 150-map feasibility figures. No dev-scale numbers remain.
% Verify each figure exists in manuscript/figures/ before final compile.

\subsection{Communication is load-bearing under sensor dropout}\label{sec:res-anchor}
We first establish that collaborative perception is not a refinement but a
requirement in this regime. Table~\ref{tab:anchor} is an
\emph{information ablation}: a single policy is evaluated on the same maps
with shared observations fused into its range channel (ON) and with its own
LiDAR only (OFF), under $\sim\!33\%$ sustained blindness at density $0.27$.
Because the weights are identical across the two arms, the gap is
attributable to the shared information alone. Sharing raises honest-drone
success from $45.9\%$ to $89.3\%$ and all-reach map success from $10.4\%$
to $67.8\%$ --- a $+57.4$ percentage-point map-level gap
(CI $[+52.8,+61.8]$). A policy \emph{trained} with sharing reaches a
similar $87.7\%$ in its native mode; we do not interpret the small
difference between the two, as it conflates two independently trained
policies. The consistent picture is that the benefit resides in the shared
information itself, which even a policy never trained with it exploits
zero-shot. A dropout sweep (Table~\ref{tab:dropout}) shows this
gap is an operating-point phenomenon: at $0\%$ dropout sharing confers no
advantage --- the gap is in fact slightly negative ($-1.3$ pp,
CI $[-2.3,-0.4]$), consistent with the early observation that cooperative
perception did not meaningfully improve over local
perception~\cite{trupercept2020} --- at $\sim\!33\%$ the gap opens to
$+41.4$ pp, and at $\sim\!50\%$ it widens further to $+50.8$ pp. We evaluate the attack and its defenses at the
$\sim\!33\%$ operating point, where collaboration is decisively valuable yet
the map remains solvable.

\begin{table}[t]\centering
\caption{Collaborative-perception anchor at density $0.27$, 500 maps.
Sharing (ON) versus own-LiDAR only (OFF) under $\sim\!33\%$ dropout.}
\label{tab:anchor}
\begin{tabular}{lccc}
\toprule
 & Drone-level & Map-level & Gap (map, 95\% CI) \\
\midrule
ON (shared)   & $89.3\%$ & $67.8\%$ & \multirow{2}{*}{$+57.4$ $[+52.8,+61.8]$} \\
OFF (own only)& $45.9\%$ & $10.4\%$ & \\
\bottomrule
\end{tabular}
\end{table}

\begin{table}[t]\centering
\caption{Dropout sweep, drone-level success, 500 maps. The comm advantage
appears only once sensing degrades.}
\label{tab:dropout}
\begin{tabular}{lcccc}
\toprule
Dropout & Blind (est.) & ON & OFF & Gap (95\% CI) \\
\midrule
$0\%$  & ${\sim}0\%$  & $88.7$ & $90.0$ & $-1.3$ $[-2.3,-0.4]$ \\
$10\%$ & ${\sim}33\%$ & $87.7$ & $46.3$ & $+41.4$ $[+38.9,+43.8]$ \\
$20\%$ & ${\sim}50\%$ & $86.2$ & $35.4$ & $+50.8$ $[+48.3,+53.3]$ \\
\bottomrule
\end{tabular}
\end{table}

\paragraph{Reading the baselines.}
Absolute success levels throughout this section should be read against the
deliberately stacked stresses. Every condition carries $\sim$33\% sustained
sensing blackout (the premise of collaboration), a calibrated obstacle
density chosen so the task is demanding but solvable, and --- as noise grows
--- position uncertainty approaching obstacle scale. The attack-free
ceiling therefore runs from $\sim$86\% at $\sigma=0$ down to $53.4\%$ at
$\sigma=0.6$; the latter is a navigation-information limit, not a defense
artifact (Section~\ref{sec:res-limit}). Crucially, the \emph{same} ceiling
underlies every arm of every comparison, so recovery and no-harm --- which
are paired differences against that common base --- remain fair measures of
the defense regardless of the absolute level.

\subsection{The minimum-fusion Byzantine vulnerability}\label{sec:res-vuln}
Because a fabricated near obstacle wins the per-ray minimum, even a small
number of traitors is damaging, and the camouflage placement is markedly
worse than the wall. Under attack with no defense, two traitors reduce
honest success from the $\sim\!86\%$ noiseless ceiling to $71.8\%$ (wall)
and $63.5\%$ (camouflage); three traitors reach $68.8\%$ and $60.9\%$
respectively. The damage \emph{saturates} between two and three traitors
(an added $\sim\!3$ pp), consistent with overlapping phantoms occupying the
same corridor. As anticipated by the fusion mechanism, the attack is
essentially independent of the dropout rate --- it overrides healthy and
degraded sensors alike --- so we hold dropout fixed and vary sensor noise
in what follows.

\subsection{Naive consistency trust is destructive under noise}\label{sec:res-naive}
Table~\ref{tab:naive} reports the naive single-frame filter against a
two-traitor wall attack as sensor noise increases. At $\sigma=0$ it is
excellent: precision $1.00$, and it recovers $+13.7$ pp. As soon as noise
appears it inverts. Detection precision collapses to $0.23$ --- more than
three quarters of its accusations are of honest neighbours --- and the
no-harm column goes deeply negative ($-27.5$ pp at $\sigma=0.4$), meaning
that \emph{with no attacker present at all}, merely enabling the filter
destroys a quarter of the swarm's success by gating out honest teammates.
Recovery under attack turns negative: the ``defense'' leaves the swarm worse
off than no defense. This is the direct empirical signature of the
$\sqrt{2}\sigma$ honest-disagreement problem of
Section~\ref{sec:methods-defense}.

\begin{table}[t]\centering
\caption{Naive single-frame trust, wall attack, $f=2$, 500 maps.
FP-harm $=$ no-harm cost; negative recovery means worse than no defense.}
\label{tab:naive}
\begin{tabular}{lccccc}
\toprule
$\sigma$ & Base & No-harm & Attack & Defense & Recovery (P/R)\\
\midrule
$0.0$ & $85.9$ & $85.9$ & $72.3$ & $85.9$ & $+13.7$ ($1.00/0.98$)\\
$0.2$ & $79.0$ & $65.1$ & $63.1$ & $62.8$ & $-0.3$ ($0.28/0.98$)\\
$0.4$ & $65.6$ & $38.1$ & $51.6$ & $38.3$ & $-13.3$ ($0.23/0.96$)\\
$0.6$ & $54.9$ & $33.3$ & $41.2$ & $33.1$ & $-8.1$ ($0.23/0.94$)\\
\bottomrule
\end{tabular}
\end{table}

\subsection{The robust filter restores safety but hits a camouflage limit}\label{sec:res-robust}
Widening the tolerance to $\varepsilon_0+4\sigma$ repairs the false
accusations: across the noise range the robust filter holds precision
$0.92$--$1.00$ and no-harm within $\pm0.3$ pp of base (the no-harm
confidence intervals span zero at every level). Against the conspicuous
wall attack it recovers gracefully. Against camouflage, however, it retains
a blind spot exactly where the paper's threat is sharpest: at $\sigma=0.6$,
$f=2$, its detection recall is only $0.13$ and its recovery is $+3.4$ pp
(CI $[+0.9,+5.8]$). The lie hidden inside the widened tolerance band cannot
be contradicted on any single frame. This is the limitation the temporal
test removes.

\subsection{Temporal offset-bias trust: the main result}\label{sec:res-temporal}
Table~\ref{tab:temporal} gives the full two-traitor results for both attack
placements across noise, comparing the single-frame robust filter with the
composed temporal filter. Three properties hold throughout. First, the
temporal filter never regresses: on the wall attack it matches or exceeds
the robust filter at every noise level. Second, its advantage opens exactly
where the robust filter fails --- at high noise under camouflage. At
$\sigma=0.6$, $f=2$, camouflage, temporal detection recall rises from
$0.13$ to $0.69$ and recovery from $+3.4$ to $+12.2$ pp
(CI $[+9.8,+14.9]$), restoring success from the attacked $37.6\%$ to
$49.8\%$ against a base of $53.4\%$. Third, this comes at no measurable cost
to honest swarms: the no-harm figure is $-0.4$ pp with a confidence
interval spanning zero.

\begin{table}[t]\centering
\caption{Temporal versus robust trust, $f=2$, 500 maps, with $95\%$ CIs on
recovery. ``nh'' is temporal no-harm. Recall shown in parentheses.}
\label{tab:temporal}
\begin{tabular}{llccccc}
\toprule
Attack & $\sigma$ & Base & Attack & Robust rec.\,(R) & Temporal rec.\,(R) & nh\\
\midrule
\multirow{4}{*}{Wall}
 & $0.0$ & $86.0$ & $71.8$ & $+13.9$ ($0.98$) & $+13.9$ ($0.98$) & $86.1$\\
 & $0.2$ & $79.1$ & $64.5$ & $+13.1$ ($0.90$) & $+14.1$ ($0.92$) & $78.8$\\
 & $0.4$ & $65.0$ & $54.8$ & $+7.8$ ($0.57$)  & $+9.2$ ($0.82$)  & $65.1$\\
 & $0.6$ & $53.4$ & $43.5$ & $+3.7$ ($0.26$)  & $+9.7$ ($0.70$)  & $53.1$\\
\midrule
\multirow{4}{*}{Camo.}
 & $0.0$ & $86.0$ & $63.5$ & $+20.3$ ($0.97$) & $+20.3$ ($0.97$) & $86.1$\\
 & $0.2$ & $79.1$ & $58.3$ & $+17.8$ ($0.86$) & $+18.6$ ($0.94$) & $78.8$\\
 & $0.4$ & $64.9$ & $47.7$ & $+11.1$ ($0.45$) & $+16.4$ ($0.81$) & $65.1$\\
 & $0.6$ & $53.4$ & $37.6$ & $+3.4$ ($0.13$)  & $\mathbf{+12.2}$ ($0.69$) & $53.0$\\
\bottomrule
\end{tabular}
\end{table}

The effect not only survives but strengthens as traitors are added, and it
does so well past the point where honest robots lose their local majority.
Table~\ref{tab:headline} isolates the worst case --- $\sigma=0.6$,
camouflage --- across traitor counts from one to seven of ten robots.
Because each honest robot is excluded from its own neighbour set, $f$
traitors leave it with at most $9-f$ honest neighbours: at $f=5$ an honest
robot's neighbourhood is already, in expectation, majority traitor, and by
$f=7$ only two of its nine potential neighbours are honest. The temporal
filter's recovery nonetheless climbs to a plateau of roughly $+15$ pp at
$f=4$--$6$ --- $+15.2$ pp (CI $[+11.9,+18.4]$) at $f=5$ and $+15.1$ pp
(CI $[+11.6,+18.6]$) at $f=6$ --- and remains $+10.9$ pp
(CI $[+7.4,+14.5]$) at $f=7$; every one of these intervals excludes zero.
The single-frame robust filter, by contrast, stays pinned near its floor and
is no longer statistically distinguishable from no defense at $f=7$
($+1.9$ pp, CI $[-1.3,+5.2]$). Temporal detection precision \emph{improves}
monotonically with the threat ($0.68\to0.82\to0.89\to0.93\to0.96\to0.97
\to0.98$): more traitors mean more genuine phantoms among the flagged links,
so the defense is most precise exactly when it matters most, while no-harm
stays flat at $\approx\!-0.4$ pp throughout.

This is the empirical content of operating without an honest local majority.
The verdict is pairwise --- each robot tests each neighbour against its own
sensing and never counts votes --- so the mechanism is structurally
independent of how many neighbours lie, and the data bear this out: the
defense's contribution is stable and significant whether honest robots are a
majority, a tie, or a minority. We are careful about what this does
\emph{not} claim. Absolute navigation success still falls as traitors are
added (undefended success drops from $44.4\%$ at $f=1$ to $\sim\!31\%$ at
$f=7$), because more liars do more damage; the temporal filter recovers a
stable, significant \emph{fraction} of that damage, it does not render
navigation traitor-count-invariant.

\begin{table}[t]\centering
\caption{Worst case ($\sigma=0.6$, camouflage) across traitor counts,
spanning the honest-majority boundary. Because each honest robot is excluded
from its own neighbour set, its neighbourhood is majority-traitor in
expectation once $f\ge5$. Recovery in pp; detection recall in parentheses;
temporal precision and no-harm in the last columns. 500 maps.}
\label{tab:headline}
\begin{tabular}{lccccc}
\toprule
$f$ & Undefended & Robust rec.\,(R) & Temporal rec.\,(R) & Prec. & No-harm\\
\midrule
$1$ & $44.4$ & $+1.9$ ($0.13$) & $+7.1$ ($0.69$)  & $0.68$ & $-0.4$\\
$2$ & $37.6$ & $+3.4$ ($0.13$) & $+12.2$ ($0.69$) & $0.82$ & $-0.4$\\
$3$ & $35.3$ & $+5.3$ ($0.12$) & $+13.6$ ($0.68$) & $0.89$ & $-0.3$\\
$4$ & $34.0$ & $+3.5$ ($0.12$) & $+14.7$ ($0.66$) & $0.93$ & $-0.3$\\
$5$ & $32.6$ & $+4.0$ ($0.13$) & $+15.2$ ($0.67$) & $0.96$ & $-0.4$\\
$6$ & $30.1$ & $+4.6$ ($0.13$) & $+15.1$ ($0.67$) & $0.97$ & $-0.4$\\
$7$ & $31.1$ & $+1.9$ ($0.13$) & $\mathbf{+10.9}$ ($0.66$) & $0.98$ & $-0.3$\\
\bottomrule
\end{tabular}
\end{table}

Figure~\ref{fig:offsethist} shows why the mechanism works: the distribution
of the aggregated offset magnitude cleanly separates honest neighbours
(concentrated near zero, as $\sqrt{2}\sigma$ noise averages out) from
fabrications (a persistent bias well above the threshold $\eta$). A probe
of this separation, before any trust gating, reports an area under the ROC
curve of $0.99$ under oracle association and $0.85$--$0.88$ under realistic
nearest-obstacle association (the latter rising with the accumulation window
to $0.88$ at the $K_{\min}=20$ operating point), both over the full 500-map
suite.

\subsection{An adaptive attacker faces a stealth/harm bind}\label{sec:res-adaptive}
A reasonable objection is that a static attacker understates the threat. We
therefore evaluate an attacker that knows the temporal filter and tunes its
fabrication to evade it, via four knobs. The clearest is the phantom
centre-offset: how far the fabricated obstacle is displaced from the real
one it hugs. Table~\ref{tab:bind} ($f=2$, $\sigma=0.6$) shows the resulting
bind. When the phantom sits on the real obstacle (offset $0$) it is
invisible to the filter (recall $0.03$) but also harmless (it blocks no new
space; harm $-0.8$ pp). As the offset grows, harm rises --- but so does
detection recall, in lockstep. The attacker cannot be both stealthy and
harmful: evading the temporal test requires collapsing the phantom onto real
geometry, which defeats the purpose of the attack.

\begin{table}[t]\centering
\caption{Stealth/harm bind: adaptive phantom centre-offset, $f=2$,
$\sigma=0.6$, 500 maps. Harm and detection recall rise together.}
\label{tab:bind}
\begin{tabular}{lcccc}
\toprule
Offset (m) & Undefended & Harm & Temporal rec. & Recall\\
\midrule
$0.0$ & $55.7$ & $-0.8$ & $+0.6$ & $0.03$\\
$0.5$ & $53.7$ & $+1.2$ & $+2.1$ & $0.07$\\
$1.0$ & $49.9$ & $+5.0$ & $+1.6$ & $0.48$\\
$1.5$ & $44.9$ & $+10.0$& $+5.6$ & $0.69$\\
$2.0$ & $42.1$ & $+12.8$& $+9.1$ & $0.70$\\
$2.5$ & $40.6$ & $+14.3$& $+8.5$ & $0.70$\\
\bottomrule
\end{tabular}
\end{table}

The bind holds across the noise range: sweeping the offset at
$\sigma\in\{0,0.2,0.4\}$ and $f\in\{1,2,3\}$ (nine configurations) yields
the same pattern in every one --- the harmless corner is undetected, the
harmful corner is caught, and detection recall at the harmful end stays high
($0.79$--$0.98$); the bind moreover persists through severe noise
($\sigma=0.6$, Table~\ref{tab:bind}), where harmful-end recall still holds at
$0.66$--$0.70$. The remaining three knobs are likewise self-defeating.
Per-frame jitter is zero-mean and so does not shift the temporal mean; it
only raises variance, and is self-defeating in practice: as the jitter grows,
the harm it inflicts falls sharply ($-6.8$ pp) while detection recall barely
moves ($0.70\to0.61$), leaving the defended success essentially unchanged.
Intermittent broadcasting (duty) dilutes the attacker's evidence but equally
dilutes its harm: at the lowest duty the harm falls to $+5.3$ pp while recall
falls to $0.41$. A surface gap large
enough to matter pushes the phantom centre outside the tolerance band, where
it is reliably caught. In no configuration does the adaptive attacker obtain
both meaningful harm and evasion.

\subsection{The navigation perception limit}\label{sec:res-limit}
Finally we distinguish the security result from an orthogonal limit. As
noise grows the attack-free base success itself falls, from $\sim\!86\%$ at
$\sigma=0$ to $53.4\%$ at $\sigma=0.6$. This is not a defense failure but a
navigation one: fine-tuning the base navigator under sensor-noise domain
randomization recovers only $\sim\!2.5$ pp at $\sigma=0.6$, indicating a
genuine perception-information limit --- position uncertainty comparable to
obstacle size cannot be undone by the policy. Crucially, the defense never
makes this worse (no-harm $\approx0$ at every level), and the two limits are
independent: the navigation limit bounds the ceiling, while the temporal
defense recovers the security gap up to that ceiling.


\section{Discussion and limitations}\label{sec:discussion}
% Source: PAPER_MASTER_PLAN sec 7 (limitations) + the mandatory disclosures
% in CLAUDE.md. Frame honestly; a reviewer must not find an undisclosed hole.

\subsection{Two independent limits}
It is important not to conflate the two limits this study exposes. The
\emph{navigation} limit (Section~\ref{sec:res-limit}) is a property of the
sensing task: at high ranging noise the position of an obstacle is
uncertain to a degree comparable to its size, and no policy can recover
information the sensor never captured; this bounds the achievable ceiling
and is unaffected by any defense. The \emph{security} limit --- the
single-frame camouflage recall collapse --- is a property of the
\emph{detector}, and it is the one our temporal rule removes. The two are
orthogonal: our contribution restores success up to the noise-limited
ceiling, not beyond it, and never depresses it.

\subsection{On the precision caveat}
The temporal detector's precision at $\sigma=0.6$ is $0.80$--$0.82$, below a
nominal $0.9$ target, which under a purely detection-theoretic reading would
be a weakness. We argue precision is the wrong figure of merit here. It was
only ever a proxy for the harm caused by false gating, and that harm is
measured directly by the no-harm column, which is statistically
indistinguishable from zero. The reconciliation is that the residual false
flags fall on ultra-stealthy camouflage buckets --- phantoms hugging real
obstacles so tightly that they barely protrude. These are simultaneously the
hardest to distinguish from honest noise \emph{and} nearly harmless, which
is exactly the stealth/harm bind seen from the defender's side: a
misclassification there costs almost nothing. Own-LiDAR fusion and the
routed heading further cushion any wrongly gated broadcast. Tightening the
threshold to raise precision sacrifices recall on genuinely harmful
phantoms, so we adopt the operating point that maximizes recovered success
rather than detector precision.

\subsection{The verifiability boundary}
Two independent observations in the literature bound what any
consistency-based defense can achieve, and both are recovered by our
results. First, trust models are blind wherever verification is impossible:
Hallyburton and Pajic report that fabrications injected into regions no
other agent observes cannot be detected by existing
methods~\cite{aerialtrust2025}. Our camouflage attack is the statistical
analogue of that spatial blind spot --- the verifier \emph{does} observe the
region, but ranging noise forces a tolerance wide enough to hide the lie
inside it. Second, an attacker that mostly reports truthfully evades
trust accumulation: TruPercept observed that unreliable agents were not
separated from honest ones by trust scores~\cite{trupercept2020}. Our
adaptive-attacker study makes this quantitative and then closes it: evasion
requires collapsing the fabrication onto real geometry, at which point it
blocks no new space. Time, rather than a stronger single-frame test, is what
converts the unverifiable into the verifiable here --- and what remains
unverifiable is, by the same bind, unharmful.

\subsection{Robustness to trust-poisoning attacks}
A recognized failure mode of consistency-based trust is that an adversary
can engineer disagreements to make the defense downweight or exclude
\emph{honest} agents~\cite{trustflip2026}.
Our evaluation confronts this directly by reporting no-harm under an
adaptive, defense-aware attacker and finding it near zero: the temporal
test does not provide a lever by which induced disagreement translates into
honest exclusion, because it thresholds a time-averaged bias that honest
neighbours, by construction, do not accumulate.

\subsection{The warm-up window}
By construction the temporal test issues no verdict for a
(receiver, neighbour, track) bucket until $K_{\min}=20$ samples have
accumulated; during those first sightings the swarm is protected only by
the single-frame robust path. In our episodes (1200 steps) this window is
negligible, but a short-lived encounter, a late-joining traitor, or a
mission with brief pairwise contact would face reduced temporal protection
--- shrinking $K_{\min}$ trades this window against verdict reliability.

\subsection{Assumptions and their consequences}
We state the modelling assumptions and their effect on the claims. (i)~The goal-direction input to the
policy is not a straight-line bearing to the goal but the gradient of a
\emph{global} shortest-path distance field, recomputed by Dijkstra's
algorithm over the full obstacle layout at the start of each episode; each
drone is thereby handed an optimal obstacle-routed heading that presupposes a
global map it could not sense on its own. We model this as an external
mission planner that supplies routed waypoints, leaving the learned policy
responsible for local control and collision avoidance --- a conventional
division of labour, but a privileged input we disclose plainly rather than
fold into the sensing model. Its effect on our claims is conservative:
because the routed heading already steers each drone around the true obstacle
field, it \emph{dampens} the influence of a fabricated obstacle, so every
attack-induced success drop we report is a lower bound on what a planner-free
agent would suffer. Retraining the policy to navigate from a raw goal
coordinate, without the routed heading, is the single largest avenue to a
stronger navigation claim and the foremost limitation of the present study. (ii)~Communication is modelled as
lossless and zero-latency within a $10$~m range; real radios drop and delay
messages, which would reduce the evidence available to the temporal test and
warrants study under realistic channels. (iii)~Trust is applied at the
\emph{neighbour} level: one sustained contradiction excludes a neighbour's
entire contribution, discarding its honest observations too. Obstacle-level
gating would be less wasteful and is natural future work. (iv)~The trust
rule is hand-designed, not learned; we regard this as a feature, since we
find no learned component is necessary for this threat class, but it is
specific to the fabricated-obstacle model studied here. (v)~The study is in
simulation, with ten drones, in two dimensions, and with circular obstacles;
the mechanism is geometry-agnostic in principle, but validation on hardware,
in three dimensions, and with richer obstacle shapes remains open.
(vi)~The temporal test assumes sensor noise is \emph{unbiased}: its
separation of honest from fabricated relies on honest offsets being
zero-mean, so a persistent per-agent measurement bias would not cancel and
could cause an honest neighbour to be flagged. A bias common to all agents
cancels harmlessly in the pairwise offset; only \emph{agent-specific}
systematic bias is a concern, and we quantify it directly by sweeping a
synthetic per-agent bias against the same noise range ($f{=}2$, camouflage,
500 maps). No-harm stays statistically indistinguishable from zero for a
miscalibration up to $0.2$~m at every noise level tested; beyond that the
own-sensing cushion is no longer sufficient, detection precision falls from
$0.9$--$1.0$ to $0.3$--$0.4$, and no-harm reaches $-11.4$ pp at bias $0.4$~m
and $-18.2$~pp at bias $0.6$~m (both at $\sigma{=}0$, the noise level at
which the fixed component of the tolerance offers the least cover). The
effect is partly self-masking at high noise --- the noise-adaptive tolerance
$\varepsilon_0+k_\sigma\sigma$ widens with $\sigma$ and incidentally absorbs
part of the bias too --- but does not vanish. Distinguishing such a bias
from fabrication is moreover possible in principle --- a miscalibrated
sensor shifts \emph{all} of a neighbour's reported obstacles by the same
amount, whereas a camouflage lie displaces only the fabricated one, so
estimating and subtracting a per-neighbour global offset before the residual
bias test would separate the two --- which we leave to future work. Learned
per-agent pose-offset correction of exactly this kind has been shown effective
against systematic bias in benign collaborative
perception~\cite{vadivelu2020}, so the residual concern is a matter of
composing that correction with our test, not an unaddressed vulnerability.


\section{Conclusion and future work}\label{sec:conclusion}
We studied Byzantine false-obstacle attacks on collaborative perception
inside a closed-loop, learned multi-robot navigation task, and traced a
complete arc from vulnerability to defense. Collaborative perception is
load-bearing under sustained sensor dropout, raising navigation success by
more than forty percentage points, but its minimum-based fusion lets a
single Byzantine teammate override honest sensing independently of the
dropout rate. The obvious countermeasure --- cross-agent consistency
checking --- is not merely weakened but \emph{destructive} under realistic
ranging noise, since honest disagreement is mistaken for lying; a
noise-aware robust variant restores safe operation but cannot detect a
camouflage fabrication hidden inside the noise band on any single frame.

Our central result is that this single-frame limit is removed in time. A
per-neighbour running mean of the offset between a reported obstacle and the
verifier's own view of it is zero-mean for honest noise and biased for a
persistent lie; thresholding this mean recovers camouflage detection recall
from $0.13$ to $0.69$ and navigation success by up to $+12.2$ percentage
points at the severest noise level, with no measurable cost to attack-free
swarms. Because the verdict is pairwise rather than a majority vote, the
recovery does not require an honest majority: it strengthens as traitors are
added, reaching roughly $+15$ percentage points at four to six traitors of
ten and remaining significant at seven, where an honest robot has only two
honest neighbours and the single-frame check is no longer distinguishable
from no defense. A
defense-aware adaptive attacker cannot escape the mechanism: evading the
temporal test forces its phantom onto real geometry, where it blocks no new
space --- a stealth/harm bind we verify at one to three traitors across
noise levels. Notably, no learned trust component is required.

Several directions follow. Removing the routed-heading dependency would let
the policy own global planning and sharpen the security claim. Applying
trust at the obstacle level rather than the neighbour level would preserve a
suspected neighbour's honest observations. Multi-frame position denoising,
which exploits the same per-frame noise independence and static-obstacle
structure that our detector relies on, is a promising route to lift the
navigation perception limit itself. Finally, validation under lossy,
delayed communication, in three dimensions, and on physical robots would
test the mechanism's assumptions in the field.


% ---------------------------------------------------------------------
\section*{CRediT authorship contribution statement}
\todo{K. Srinivasa Rao: Conceptualization, Methodology, Software,
Investigation, Formal analysis, Writing -- original draft.
[Professor]: Supervision, Writing -- review \& editing, Project
administration.}

\section*{Declaration of competing interest}
The authors declare that they have no known competing financial interests or
personal relationships that could have appeared to influence the work
reported in this paper.

\section*{Data availability}
The simulation environment, trained models, evaluation scripts, and raw
result logs that reproduce every table and figure in this article are
publicly available at \todo{GitHub URL} (archived at \todo{Zenodo DOI}).

\section*{Declaration of generative AI and AI-assisted technologies in the
manuscript preparation process}
During the preparation of this work the authors used Claude (Anthropic) in
order to assist with drafting and editing the manuscript text. After using
this tool, the authors reviewed and edited the content as needed and take
full responsibility for the content of the published article.

\section*{Acknowledgements}
\todo{Optional: institute compute resources, colleagues, etc.}

% ---------------------------------------------------------------------
\bibliographystyle{elsarticle-num}
\bibliography{refs}

\end{document}
